% docs/texte/00_frontmatter/environment.tex -- chapter "System Environment" (Interferosim)
% Version: 0.2.0 (2026-10-12) -- project virtual environment
%   0.1.0: placeholder
% ================================================================

Python~$\geq$~3.11 with numpy, scipy, matplotlib and astropy; JupyterLab,
ipywidgets and ipykernel for the notebooks; pytest and nbval for the
tests. The notebooks run locally, on JupyterHub or on Binder
(\texttt{binder/environment.yml}).

\section{Project virtual environment}
\index{virtual environment}

All work uses the virtual environment \texttt{.venv} in the project
folder (excluded from version control by \texttt{.gitignore}). It is
created by a script:

\begin{center}
\begin{tabular}{@{}ll@{}}
\toprule
System & Command (in the project folder) \\
\midrule
Windows & \texttt{setup\_venv.bat} (double-click) or\\
        & \texttt{powershell -ExecutionPolicy Bypass -File setup\_venv.ps1} \\
Linux, macOS & \texttt{./setup\_venv.sh} \\
\bottomrule
\end{tabular}
\end{center}

The script searches for Python~$\geq$~3.11, creates \texttt{.venv}, installs
Interferosim in editable mode with the extras \texttt{notebooks} and
\texttt{test}, registers the Jupyter kernel \enquote{Python
(interferosim)} and runs the tests. Options of the Windows script:
\texttt{-Python <path>} (explicit interpreter), \texttt{-Recreate}
(delete and rebuild \texttt{.venv}), \texttt{-NoTests},
\texttt{-DropboxIgnore} (exclude \texttt{.venv} from Dropbox
synchronisation).

After activation (\texttt{.venv\textbackslash Scripts\textbackslash activate}
on Windows, \texttt{source .venv/bin/activate} on Linux/macOS) the command
line tool \texttt{interferosim-sky} and \texttt{jupyter lab} are available.
